% Introduction Slides

\begin{frame}{Bilinear MLPs: Weight-Based Interpretability}
    \begin{columns}
        \column{0.5\textwidth}
        \textbf{Purpose:}
        \begin{itemize}
            \item Direct interpretability from weights
            \item No post-hoc analysis needed
            \item No SAE training required
        \end{itemize}
        
        \vspace{0.3cm}
        \textbf{Key Insight:}\\
        Bilinear structure creates \textbf{quadratic} interactions that can be \textbf{eigendecomposed} to reveal interpretable input directions.
        
        \column{0.5\textwidth}
        \textbf{Method:}
        \begin{itemize}
            \item Replace nonlinearity with element-wise multiplication
            \item This creates an \textbf{interaction matrix} for each output
            \item Eigendecompose to find dominant input directions
            \item Top eigenvectors = most important patterns
        \end{itemize}
    \end{columns}
\end{frame}

\begin{frame}{Mathematical Foundation}
    \small
    \textbf{Bilinear Layer:} $\mathbf{y} = (\mathbf{W}_l \mathbf{x}) \odot (\mathbf{W}_r \mathbf{x})$
    
    \vspace{0.15cm}
    \textbf{Output as Quadratic Form:} For each output $c$: \quad $y_c = \mathbf{x}^\top B_c \, \mathbf{x}$
    
    \vspace{0.05cm}
    where $B_c$ is the \textbf{interaction matrix}---captures how input pairs $(x_i, x_j)$ contribute to $y_c$.
    
    \vspace{0.15cm}
    \textbf{Eigendecomposition:} Symmetrize and decompose $B_c$:
    \begin{equation*}
        B_c = \sum_{r=1}^{d} \lambda_r^{(c)} \mathbf{v}_r^{(c)} (\mathbf{v}_r^{(c)})^\top
        \quad \Rightarrow \quad
        y_c = \sum_r \lambda_r (\mathbf{v}_r^\top \mathbf{x})^2
    \end{equation*}
    
    \vspace{-0.1cm}
    \footnotesize
    Eigenvectors $\mathbf{v}_r^{(c)} \in \mathbb{R}^{d_{\text{in}}}$ = interpretable input directions; eigenvalues $\lambda_r^{(c)}$ = importance.
    
    \vspace{0.15cm}
    \small
    \textbf{Effective Rank:} Measures spectral concentration (how many eigenvalues matter):
    \begin{equation*}
        \text{EffRank}(\boldsymbol{\lambda}) = \left( \frac{\|\boldsymbol{\lambda}\|_1}{\|\boldsymbol{\lambda}\|_2} \right)^2
    \end{equation*}
    
    \vspace{-0.15cm}
    \footnotesize
    Low effective rank $\Rightarrow$ few dominant eigenvectors $\Rightarrow$ more interpretable.
\end{frame}

\begin{frame}{Two Analysis Methods (Paper Section 3)}
    \textbf{The analysis method depends on what features are available:}
    
    \vspace{0.15cm}
    \begin{columns}
        \column{0.5\textwidth}
        \fcolorbox{blue!50}{blue!10}{
            \parbox{0.95\textwidth}{
                \small
                \textbf{Method A: No Prior Features}\\[0.1cm]
                (Vision, Section 4)
                
                \vspace{0.1cm}
                \textbf{Model:} Single bilinear MLP + linear head (MNIST/Fashion-MNIST).
                
                \vspace{0.1cm}
                Classes act as output features.\\
                Eigendecompose $B_c$ directly:
                \begin{equation*}
                    y_c = \sum_r \lambda_r (\mathbf{v}_r^\top \mathbf{x})^2
                \end{equation*}
                Eigenvectors $\mathbf{v}_r$ visualized as 28$\times$28 images.
            }
        }
        
        \column{0.5\textwidth}
        \fcolorbox{green!50}{green!10}{
            \parbox{0.95\textwidth}{
                \small
                \textbf{Method B: SAE Features}\\[0.1cm]
                (Language, Section 5)
                
                \vspace{0.1cm}
                \textbf{Model:} Transformer with bilinear MLPs (replacing standard MLPs) + pretrained SAEs.
                
                \vspace{0.1cm}
                SAEs $\to$ sparse features $\mathbf{z}^{\text{in}}, \mathbf{z}^{\text{out}}$.\\
                Build $Q_f$ per output feature $f$:
                \begin{equation*}
                    z^{\text{out}}_f \approx (\mathbf{z}^{\text{in}})^\top Q_f \, \mathbf{z}^{\text{in}}
                \end{equation*}
            }
        }
    \end{columns}
    
    \vspace{0.2cm}
    \centering
    \small
    \textbf{Both methods} use eigendecomposition to identify important input directions for each output.
\end{frame}
